% GENERATED FILE. Edit canonical lessons, then run npm run generate:books.
% canonical-source-hash: fnv1a64:646c0f9dd76cfece
% canonical-lessons: PA-W09-date-label, PA-W09-month-map, PA-W09-date-format, PA-W09-date-a, PA-W09-date-b, PA-W09-date-select, PA-W09-date-supported, PA-W09-date-delayed, PA-W09-date-dictation, PA-W09-date-repair, PA-W09-date-no-model

\chapter{Write One Fictional Date Field}
\label{ch:pa-form-date-runway}
\hlchaptermodality{153}

\begin{chapteropening}
\textbf{By the end of this chapter:} \emph{I can recognise a Punjabi date label, connect two taught month names to their numeric positions, and independently write one of two fictional dates in day/month/year order from a nonverbal cue.}

The learner fills one fictional date line from a remembered nonverbal cue without a value bank, romanization, Latin digits, or a copyable Gurmukhi answer, then checks cue selection, month, digit order, separators, and placement separately.
\end{chapteropening}

\section[\texorpdfstring{\pa{ਤਾਰੀਖ਼} (tārīkh)}{ਤਾਰੀਖ਼ (tārīkh)}]{\pa{ਤਾਰੀਖ਼} — find the date line}

\label{lesson:PA-W09-date-label}



\begin{quote}
You met \textbf{\pa{ਤਾਰੀਖ਼}} as “a date.” Read it aloud once. Its Persian source \emph{tārīkh} also means a date or history; here the form asks for a date only.
\end{quote}

\begin{tcolorbox}[breakable,skin=enhanced,colback=orange!6,colframe=orange!45!black,boxrule=0.5pt,arc=1mm,left=6pt,right=6pt,top=4pt,bottom=4pt,fonttitle=\bfseries,title={Read this}]
On this fictional beginner form, \textbf{\pa{ਤਾਰੀਖ਼}: \_\_\_\_\_\_} labels the date line. Point to the label. The blank is where a date goes; it is not a phone or age line. No real personal information is requested.
\end{tcolorbox}

\subsection*{Your turn — copy the label}
Copy \textbf{\pa{ਤਾਰੀਖ਼}:} once, including the colon. Leave its value blank.

\begin{tcolorbox}[breakable,colback=teal!4,colframe=teal!35!black,title={Before you move on}]
Next, we will connect two month names already met to their small numeric positions. There is no full date to write yet.
\end{tcolorbox}
\section[\texorpdfstring{\pa{ਜਨਵਰੀ} · \pa{ਫ਼ਰਵਰੀ} (janvarī · farvarī)}{ਜਨਵਰੀ · ਫ਼ਰਵਰੀ (janvarī · farvarī)}]{Two known months, two short codes}

\label{lesson:PA-W09-month-map}



\begin{quote}
Say the two month names already met: \textbf{\pa{ਜਨਵਰੀ}} (January), then \textbf{\pa{ਫ਼ਰਵਰੀ}} (February). Do not try to memorise a new list of months.
\end{quote}

\begin{grammarlens}[title={Grammar lens — month position}]
In a two-place month box, January is \textbf{\pa{੦੧}} and February is \textbf{\pa{੦੨}}. The first place is the already-taught zero; the second is one or two. Point to each Gurmukhi digit while saying the month name. These are month \emph{positions}, not new ways to pronounce the words.
\end{grammarlens}

\subsection*{Your turn — one place at a time}
Copy \textbf{\pa{੦੧}} under \pa{ਜਨਵਰੀ} and \textbf{\pa{੦੨}} under \pa{ਫ਼ਰਵਰੀ}. Cover the words; when you hear “February,” uncover and point to \textbf{\pa{੦੨}}. Repair only the differing digit.

\begin{tcolorbox}[breakable,colback=teal!4,colframe=teal!35!black,title={Before you move on}]
Only these two month positions are in this closed practice bank.
\end{tcolorbox}
\section[\_\_/\_\_/\_\_\_\_ (day/month/year)]{Three small boxes, two slashes}

\label{lesson:PA-W09-date-format}



\begin{quote}
Point to \textbf{\pa{੦੧}} for January and \textbf{\pa{੦੨}} for February. We will put one of these in the \emph{middle} of a date, not at the beginning.
\end{quote}

\begin{grammarlens}[title={Grammar lens — the form's printed order}]
This fictional form explicitly prints \textbf{day / month / year}. It has two places for the day, two for the month, and four for the year: \textbf{\_\_/\_\_/\_\_\_\_}. Copy the two slash marks between the boxes. A slash is a separator, never a digit. Other forms may use another order; follow the printed order on this one.
\end{grammarlens}

\subsection*{Your turn — locate the month box}
Fill only the middle box with \textbf{\pa{੦੧}}. Leave the day and year boxes empty: \textbf{\_\_/\pa{੦੧}/\_\_\_\_}. Point to the two slash marks. Then cover this model and draw two empty slash-separated boxes again.

\begin{tcolorbox}[breakable,colback=teal!4,colframe=teal!35!black,title={Before you move on}]
The next lessons build two complete fictional examples, one piece at a time.
\end{tcolorbox}
\section[\texorpdfstring{\pa{੧੫}/\pa{੦੧}/\pa{੨੦੨੫} (fictional January date A)}{੧੫/੦੧/੨੦੨੫ (fictional January date A)}]{Fictional date A — one box at a time}

\label{lesson:PA-W09-date-a}



\begin{quote}
Trace the three empty boxes with a finger: day / month / year. The slash marks stay in place.
\end{quote}

\begin{tcolorbox}[breakable,skin=enhanced,colback=orange!6,colframe=orange!45!black,boxrule=0.5pt,arc=1mm,left=6pt,right=6pt,top=4pt,bottom=4pt,fonttitle=\bfseries,title={Read this}]
Practice card A says “the fifteenth of January, two thousand twenty-five.” It is fictional. The day is \textbf{\pa{੧੫}}, January is \textbf{\pa{੦੧}}, and the four year digits are \textbf{\pa{੨੦੨੫}}. Put the pieces together: \textbf{\pa{੧੫}/\pa{੦੧}/\pa{੨੦੨੫}}. Touch and read only one piece at a time; no real birthday is requested.
\end{tcolorbox}

\subsection*{Your turn — copy in three passes}
Copy the day \textbf{\pa{੧੫}}, add a slash, copy month \textbf{\pa{੦੧}}, add a slash, then copy the year \textbf{\pa{੨੦੨੫}}. Check that the two slash marks divide two, two, and four digits.

\begin{tcolorbox}[breakable,colback=teal!4,colframe=teal!35!black,title={Before you move on}]
Keep A visible for now. Independent recall comes later.
\end{tcolorbox}
\section[\texorpdfstring{\pa{੨੫}/\pa{੦੨}/\pa{੨੦੨੫} (fictional February date B)}{੨੫/੦੨/੨੦੨੫ (fictional February date B)}]{Fictional date B — change two small pieces}

\label{lesson:PA-W09-date-b}



\begin{quote}
Point to A, \textbf{\pa{੧੫}/\pa{੦੧}/\pa{੨੦੨੫}}. Its middle box means January.
\end{quote}

\begin{tcolorbox}[breakable,skin=enhanced,colback=orange!6,colframe=orange!45!black,boxrule=0.5pt,arc=1mm,left=6pt,right=6pt,top=4pt,bottom=4pt,fonttitle=\bfseries,title={Read this}]
Practice card B says “the twenty-fifth of February, two thousand twenty-five.” The day changes to \textbf{\pa{੨੫}}; February changes the middle box to \textbf{\pa{੦੨}}; the year stays \textbf{\pa{੨੦੨੫}}. The complete fictional date is \textbf{\pa{੨੫}/\pa{੦੨}/\pa{੨੦੨੫}}. Compare one box at a time with A.
\end{tcolorbox}

\subsection*{Your turn — copy, then compare}
Copy B once. Circle only the two boxes that differ from A. Say “day” for the first and “month” for the second; do not change the year or slash marks.

\begin{tcolorbox}[breakable,colback=teal!4,colframe=teal!35!black,title={Before you move on}]
These two examples form a deliberately closed practice bank, not an invitation to write a real personal date.
\end{tcolorbox}
\section[\texorpdfstring{\pa{ਕ} · \pa{ਖ} (two nonverbal selectors)}{ਕ · ਖ (two nonverbal selectors)}]{Select the requested fictional date}

\label{lesson:PA-W09-date-select}



\begin{quote}
Read the closed bank once: A \textbf{\pa{੧੫}/\pa{੦੧}/\pa{੨੦੨੫}}; B \textbf{\pa{੨੫}/\pa{੦੨}/\pa{੨੦੨੫}}. Point to the middle box in each. A is January; B is February.
\end{quote}

\subsection*{You'll want to know — a symbol, not a copied answer}
The fictional card marked \textbf{\pa{ਕ}} requests A. The card marked \textbf{\pa{ਖ}} requests B. These single-letter shapes are selectors, not month names or dates. The same two shapes were used in the earlier bounded form-field practice.

\subsection*{Your turn — decide before writing}
Point to \textbf{\pa{ਖ}}. Say “B” before you look at the date. Then point to B in the bank; only then copy it. Repeat with \textbf{\pa{ਕ}} and A. Separate \emph{selection} from digit writing.

\begin{tcolorbox}[breakable,colback=teal!4,colframe=teal!35!black,title={Before you move on}]
The bank remains visible here. It will be hidden only after supported practice.
\end{tcolorbox}
\section[\texorpdfstring{\pa{ਤਾਰੀਖ਼}: \pa{੧੫}/\pa{੦੧}/\pa{੨੦੨੫} (supported …)}{ਤਾਰੀਖ਼: ੧੫/੦੧/੨੦੨੫ (supported …)}]{Fill one date line with help}

\label{lesson:PA-W09-date-supported}



\begin{quote}
Point to the \textbf{\pa{ਤਾਰੀਖ਼}} label. Keep the closed bank in view: A \textbf{\pa{੧੫}/\pa{੦੧}/\pa{੨੦੨੫}}, B \textbf{\pa{੨੫}/\pa{੦੨}/\pa{੨੦੨੫}}.
\end{quote}

\subsection*{Writing — supported entry}
For the fictional card \textbf{\pa{ਕ}}, select A. Copy its day into the first box, January's numeric month into the middle box, and its year into the final box: \textbf{\pa{ਤਾਰੀਖ਼}: \pa{੧੫}/\pa{੦੧}/\pa{੨੦੨੫}}. Keep both slash marks.

\subsection*{Your turn — check one dimension}
Read the line from left to right. Check cue, day, month, year, and separators one at a time. If a piece differs, change only that piece.

\begin{tcolorbox}[breakable,colback=teal!4,colframe=teal!35!black,title={Before you move on}]
This is supported copying, not the independent checkpoint.
\end{tcolorbox}
\section[\texorpdfstring{\pa{ਤਾਰੀਖ਼}: \pa{੨੫}/\pa{੦੨}/\pa{੨੦੨੫} (delayed …)}{ਤਾਰੀਖ਼: ੨੫/੦੨/੨੦੨੫ (delayed …)}]{Hide the bank, then write}

\label{lesson:PA-W09-date-delayed}



\begin{quote}
Study B \textbf{\pa{੨੫}/\pa{੦੨}/\pa{੨੦੨੫}} for five seconds. Touch day, month, and year, then cover this line and the earlier date lessons.
\end{quote}

\subsection*{Writing — delayed entry}
After the short delay, the fictional card shows \textbf{\pa{ਖ}}. Fill its \textbf{\pa{ਤਾਰੀਖ਼}} line from memory. Do not uncover the bank before you finish. Write the day, the month, and the year in the printed order, with two slash marks.

\subsection*{Your turn — uncover and repair}
Uncover B only after the attempt. Compare one box at a time. If the month is wrong, correct only the middle two digits; keep the other boxes unchanged.

\begin{tcolorbox}[breakable,colback=teal!4,colframe=teal!35!black,title={Before you move on}]
This is delayed practice. The final no-model checkpoint still comes later.
\end{tcolorbox}
\section[\texorpdfstring{\pa{੧੫}/\pa{੦੧}/\pa{੨੦੨੫} (heard fictional date A)}{੧੫/੦੧/੨੦੨੫ (heard fictional date A)}]{Hear a date; fill three boxes}

\label{lesson:PA-W09-date-dictation}



\begin{quote}
Draw three empty boxes separated by slashes. Say “day, month, year” while pointing. Cover all written date examples.
\end{quote}

\subsection*{You'll want to know — one short fictional prompt}
Have a partner or the narration read: “the fifteenth of January, two thousand twenty-five.” Listen once; then listen again. Do not display a digit model while you write.

\subsection*{Writing — transcription}
Write only the three heard pieces in Gurmukhi digits. Put the day in the first box, January's taught two-digit month in the middle, and the year last. Then check only the two slash positions.

\begin{tcolorbox}[breakable,colback=teal!4,colframe=teal!35!black,title={Before you move on}]
No real birth date or other personal data is being collected.
\end{tcolorbox}
\section[\texorpdfstring{\pa{ਤਾਰੀਖ਼}: \_\_/\_\_/\_\_\_\_ (date field repair)}{ਤਾਰੀਖ਼: \_\_/\_\_/\_\_\_\_ (date field repair)}]{Repair one dimension, not the whole line}

\label{lesson:PA-W09-date-repair}



\begin{quote}
Read A and B once. A is January; B is February. Their year is the same.
\end{quote}

\begin{grammarlens}[title={Grammar lens — a short repair order}]
Check one dimension at a time: \textbf{selector $\to$ day $\to$ month $\to$ year $\to$ slash marks $\to$ label spelling and spacing}. On a printed date form, a correct set of digits in the wrong boxes is still wrong. The label is \textbf{\pa{ਤਾਰੀਖ਼}}; keep the final dot beneath \textbf{\pa{ਖ਼}}, and leave a readable space after the colon.
\end{grammarlens}

\subsection*{Your turn — two tiny repairs}
For \textbf{\pa{ਕ}}, \textbf{\pa{੧੫}/\pa{੦੧੨੦੨੫}} has the right digits but one missing slash. Add only that slash. Next, \textbf{\pa{ਤਾਰੀਖ}: \pa{੧੫}/\pa{੦੧}/\pa{੨੦੨੫}} has the right value but lacks the dot in the label; repair only \textbf{\pa{ਤਾਰੀਖ਼}}. Keep the date digits untouched.

\begin{tcolorbox}[breakable,colback=teal!4,colframe=teal!35!black,title={Before you move on}]
If the value is right but the label or spacing is wrong, repair that later as a separate pass. Do not erase a correct date to fix one mark.
\end{tcolorbox}
\section[\texorpdfstring{\pa{ਤਾਰੀਖ਼}: \_\_\_\_\_\_\_\_\_\_ (independent …)}{ਤਾਰੀਖ਼: \_\_\_\_\_\_\_\_\_\_ (independent …)}]{One no-model fictional date field}

\label{lesson:PA-W09-date-no-model}



\begin{quote}
Close the earlier lessons. There is no value bank, support-language label, Latin-digit version, or copyable Gurmukhi date below.
\end{quote}

\subsection*{Writing — independent controlled choice}
Complete the requested fictional field:

\begin{quote}
\pa{ਖ} — \textbf{\pa{ਤਾਰੀਖ਼}: \_\_\_\_\_\_\_\_\_\_}
\end{quote}

Choose the known date attached to the selector. Write its day, month, and year from memory in Gurmukhi digits with the two printed separators. Do not reopen a lesson until the attempt is complete.

\subsection*{Your turn — separate repair passes}
After writing, compare selector, day, month, year, separators, label spelling, and spacing separately. Repair only the first differing dimension.

\begin{tcolorbox}[breakable,colback=teal!4,colframe=teal!35!black,title={Before you move on}]
This closes only the untimed date-field ladder. Integration with other form fields, timed writing, two full mocks, and human A1 validation remain separate dependency-linked work.
\end{tcolorbox}
